\documentclass[journal,twoside,web]{ieeecomm}
\usepackage{cite}
\usepackage{amsmath,amssymb,amsfonts,amsthm}
\usepackage{graphicx}
\usepackage{algorithm,algorithmic}
\usepackage[caption=false]{subfig}
\usepackage{tikz}
\usepackage{hyperref}
\hypersetup{hidelinks}
% Journal name for the running head and the accent colour (normally supplied
% by a journal .sty file such as generic.sty or tac.sty).
\def\journalname{IEEE TRANSACTIONS AND JOURNALS TEMPLATE}
\definecolor{subsectioncolor}{rgb}{0,0.263,0.576}
\newtheorem{theorem}{Theorem}
\begin{document}
\title{Sample Article for the Rebuilt Colour Class with $\mathcal{H}_\infty$ and \revised{Highlighted} Maths}
\author{First A. Author, \IEEEmembership{Fellow, IEEE}, Second B. Author, and Third C. Author Jr., \IEEEmembership{Member, IEEE}
\thanks{This paragraph of the first footnote has a blank line in it.

Second paragraph of the same \texttt{\string\thanks}, which fails in V1.7.}
\thanks{Second B. Author Jr. was with Rice University, Houston, TX 77005 USA.}}
\markboth{Journal of \LaTeX\ Class Files, Vol.~14, No.~8, August 2021}%
{Author \MakeLowercase{\textit{et al.}}: A Test of the Rebuilt Colour Class}
\maketitle
\begin{abstract}
This abstract contains inline maths $x^2 + \mu_0 H$ that must be in the normal math italic, not upright sans. \revised{This sentence is a revision inside the abstract.}
\end{abstract}
\begin{IEEEkeywords}
Index terms, $\ell_1$ norm, testing.
\end{IEEEkeywords}

\section{Introduction}
\IEEEPARstart{T}{his} document exercises the class. A citation \cite{a,b,c} and a reference to Section~\ref{sec:two}, Fig.~\ref{fig:one}, Table~\ref{tab:one} and \eqref{eq:one}.

\revised{This whole paragraph is a revision. It has inline maths $a = b + \frac{1}{2}$, and it continues.

It also has a second paragraph, and a display:
\begin{equation}
  \int_0^\infty e^{-st} f(t)\,dt = F(s), \label{eq:one}
\end{equation}
and text after the display.} Text after the revision is black again, with maths $c = d$.

Inline revised maths: $a + \revised{b} = c$, $x \revised{=} y$, and $\revised{\sum_{k=1}^n k}$.

\begin{revision}
A block revision with an alignment:
\begin{align}
  a &= b + c, \label{eq:al}\\
  d &= e - f.
\end{align}
And an IEEEeqnarray:
\begin{IEEEeqnarray}{rCl}
  g & = & h \IEEEyesnumber \IEEEyessubnumber \\
  i & = & j. \IEEEyessubnumber
\end{IEEEeqnarray}
\begin{itemize}
\item A list item, revised.
\item Another.
\end{itemize}
\end{revision}
Black again.

\subsection{Subsection Heading With Maths $\mathbf{H}$}
Text. \revised{Short.} More.
\subsubsection{Run-in Heading $x$}
Text following the run-in heading continues in the same paragraph and must be black and 10pt.
\paragraph{Paragraph Heading}
Text.

\section*{Unnumbered Section}
Must be centred exactly like a numbered one.

\begin{figure}[!t]
\centering
\begin{tikzpicture}\fill[nblue] (0,0) rectangle (2,1); \draw[subsectioncolor] (0,0) -- (2,1);\end{tikzpicture}
\caption{Short caption, centred.}
\label{fig:one}
\end{figure}

\begin{figure}[!t]
\centering
\subfloat[First]{\rule{1in}{0.75in}\label{fig:a}}
\hfil
\subfloat[Second]{\rule{1in}{0.75in}\label{fig:b}}
\caption{A longer caption that certainly wraps around onto a second line so that the wrapped branch of the caption code is exercised, with \revised{a revised phrase} and maths $\alpha$.}
\end{figure}

\begin{table}[!t]
\caption{A Table Caption in Small Caps}
\label{tab:one}
\centering
\begin{tabular}{c|c}
\hline
One & Two\\
\hline
$x$ & $y$\\
\hline
\end{tabular}
\end{table}

\begin{table*}[!t]
\caption{A Double Column Table}
\centering
\begin{tabular}{c|c}
\hline
One & Two\\
\hline
\end{tabular}
\end{table*}

\begin{algorithm}[!t]
\caption{An algorithm}
\begin{algorithmic}
\STATE $x \leftarrow 0$
\end{algorithmic}
\end{algorithm}

\section{Second}
\label{sec:two}
\begin{theorem}
An amsthm theorem.
\end{theorem}
\begin{proof}
An amsthm proof, since amsthm is loaded.
\end{proof}
\begin{IEEEproof}
An IEEEtran proof.
\end{IEEEproof}
Filler. Filler text. Filler text. Filler text. Filler text. Filler text. Filler text. Filler text. Filler text. Filler text. Filler text. Filler text.

\appendices
\section{First Appendix}
Appendix text.

\section*{Acknowledgment}
Thanks.

\begin{thebibliography}{1}
\bibitem{a} A. Author, ``A paper,'' \emph{IEEE Trans. Autom. Control}, vol. 1, no. 1, pp. 1--2, 2020.
\bibitem{b} B. Author, ``Another paper,'' \emph{IEEE Trans. Autom. Control}, vol. 1, no. 1, pp. 1--2, 2020.
\bibitem{c} C. Author, ``A third paper,'' \emph{IEEE Trans. Autom. Control}, vol. 1, no. 1, pp. 1--2, 2020.
\end{thebibliography}

\begin{IEEEbiography}[{\rule{1in}{1.25in}}]{First A. Author}
Biography text with maths $\sigma$. Biography text. Biography text. Biography text. Biography text. Biography text.
\end{IEEEbiography}
\begin{IEEEbiographynophoto}{Second B. Author}
No photo biography.
\end{IEEEbiographynophoto}
\begin{biography}{Third C. Author}
Old-name biography.
\end{biography}
\end{document}
